\documentclass[10pt,a4paper]{amsart}
\usepackage[margin=19mm]{geometry}
\usepackage{amsmath,amssymb,mathtools}
\usepackage[scr=rsfs,cal=euler]{mathalfa}
\usepackage{xfrac}
\usepackage[unicode,hidelinks,bookmarks=false]{hyperref}
\newcommand{\ie}{i.e.\ }
\newcommand{\cf}{cf.\ }
\newcommand{\resp}{respectively\ }
\newcommand{\Subsection}{\subsection}
\providecommand{\nobreakdash}{\nobreak\hbox{-}}
\setlength{\emergencystretch}{2em}
\multlinegap0pt
\usepackage{mathtools}
\usepackage{amssymb}
\usepackage[shortlabels]{enumitem}
\usepackage{tikz}
\usepackage{cancel}
\newtheorem{theorem}{Theorem}
\title{Rigidity and Positivity of Hawking Quasi-Local Energy on Area-Constrained Critical Surfaces}
\author{Alejandro Pe\~nuela Diaz}
\date{Source: arXiv:2507.16588v3 (CC BY 4.0)}
\begin{document}
\maketitle

\noindent\emph{Source and licence.} Rigidity and Positivity of Hawking Quasi-Local Energy on Area-Constrained Critical Surfaces. Version arXiv:2507.16588v3 (CC BY 4.0), distributed by arXiv under the Creative Commons Attribution 4.0 International licence, http://creativecommons.org/licenses/by/4.0/, as recorded in the arXiv licence metadata for this exact version and shown on the abstract page https://arxiv.org/abs/2507.16588v3. Full licence notice: \texttt{licenses/arxiv-2507.16588-CC-BY-4.0.txt}.\par\smallskip

\noindent\textit{Editorial scope.} Counted unit: the article's theorem thm:unified\_positivity (source0.tex lines 688--699). The article's complete original proof of it is delivered with it (source0.tex lines 700--725, one \texttt{proof} environment of 26 lines). No mathematics is written, repaired or re-derived: the statement and the proof are the article's own lines, in the article's own notation, and no retained line is substituted or re-flowed.  Retained uncounted as the source-local context the proof needs: lines 667--669.  Nothing was omitted from any retained span, so the package carries no omission record.  No retained line contains a citation, so the wrapper carries no bibliography and no bibliography entry.  No retained line contains a figure environment, an image-inclusion command or drawing code, so no image is delivered and the package declares no asset dependency.  Editorial formatting only, in addition: an AMS \texttt{amsart} preamble that restates the article's own declarations and macros used by the retained lines, namely the environments \texttt{theorem} on separate counters, together with the standard packages those lines need.  The retained environments print in this document as Theorem 1 in order of appearance, whereas the article prints the same items with the numbering of its own full document.  The complete native source of the article as distributed in the arXiv e-print for this exact version is bundled unchanged as \texttt{sources/182021/source0.tex}.
\begin{equation}\label{willmorehigherdim}
    0= \lambda H  +\Delta^\Sigma H- \frac{n-3}{2(n-1)}H^3 + H|\mathring{B}|^2+ H \mathrm{Ric}^M(\nu, \nu) 
\end{equation}

\begin{theorem}\label{thm:unified_positivity}
Let $(M,g)$ be an $n$-dimensional Riemannian manifold ($n\geq 3$), and let $\Lambda \in \mathbb{R}$ be a constant. Let $\Sigma \subset M$ be a closed, area-constrained Willmore hypersurface with strictly positive mean curvature ($H > 0$) and parameter $\lambda$. Assume that the ambient scalar curvature satisfies $\mathrm{Sc}^M \ge 2\Lambda$ along $\Sigma$.
\begin{itemize}
    \item[$(i)$] If $\label{unified_bound_1}
    \lambda \geq \frac{n-3}{2(n-2)|\Sigma|} \int_{\Sigma} \mathrm{Sc}^{\Sigma} d\mu - \frac{n-1}{n}\Lambda$,
then $\mathcal{E}_{n,\Lambda,1}(\Sigma)\geq 0$.
    \item[$(ii)$] If $
    \lambda \geq \frac{1}{2|\Sigma|}\int_{\Sigma} \mathrm{Sc}^{\Sigma} d\mu - \frac{n-1}{2} \left( \frac{\omega_{n-1}}{ |\Sigma|} \right)^{\frac{2}{n-1}} - \frac{n-1}{n}\Lambda$,
then $\mathcal{E}_{n,\Lambda,2}(\Sigma)\geq 0$.
\end{itemize}
Furthermore, if either inequality is strict, the respective Hawking energy is strictly positive.
\end{theorem}

\begin{proof}
Because $H > 0$ strictly, we may multiply the Willmore equation \eqref{willmorehigherdim} by $H^{-1}$. Integrating the Laplacian term by parts and using the Gauss equation to substitute $\mathrm{Ric}^M(\nu, \nu)$ yields the fundamental integral identity:
\begin{equation}\label{willmoredived_unified}
\lambda |\Sigma | + \int_{\Sigma} |\nabla^{\Sigma} \log H|^2 + \frac{1}{2(n-1)} H^2 + \frac{1}{2} |\mathring{B}|^2 \, d\mu = \frac{1}{2} \int_{\Sigma } (\mathrm{Sc}^{\Sigma} - \mathrm{Sc}^M) \, d\mu.
\end{equation}
Applying the ambient scalar curvature bound $\mathrm{Sc}^M \ge 2\Lambda$, the right-hand side is strictly bounded above by $\frac{1}{2} \int_{\Sigma } \mathrm{Sc}^{\Sigma} d\mu - \Lambda|\Sigma|$. 

Assuming the first condition  for $\lambda$, we obtain
\begin{equation*}
\begin{split}
    \int_{\Sigma} \frac{1}{2(n-1)} H^2 \, d\mu &\leq \frac{1}{2} \int_{\Sigma} \mathrm{Sc}^{\Sigma} d\mu - \Lambda|\Sigma| - \lambda|\Sigma| \\
    &\leq \frac{n-2 - (n-3)}{2(n-2)} \int_{\Sigma} \mathrm{Sc}^{\Sigma} d\mu - \left(1 - \frac{n-1}{n}\right)\Lambda|\Sigma| \\
    &= \frac{1}{2(n-2)} \int_{\Sigma} \mathrm{Sc}^{\Sigma} d\mu - \frac{1}{n}\Lambda|\Sigma|.
\end{split}
\end{equation*}
Multiplying this inequality algebraically by $2(n-1)(n-2)/|\Sigma|$  implies $\mathcal{E}_{n,\Lambda,1}(\Sigma)\geq 0$.

If instead we assume the second condition  for $\lambda$, the same substitution yields
\begin{equation*}
\begin{split}
    \int_{\Sigma} \frac{1}{2(n-1)} H^2 \, d\mu &\leq \frac{1}{2} \int_{\Sigma} \mathrm{Sc}^{\Sigma} d\mu - \Lambda|\Sigma| - \lambda|\Sigma| \\
    &\leq \frac{n-1}{2} \left( \frac{\omega_{n-1}}{ |\Sigma|} \right)^{\frac{2}{n-1}}|\Sigma| - \frac{1}{n}\Lambda|\Sigma|.
\end{split}
\end{equation*}
Dividing this by $\frac{n-1}{2}$  recovers the bound $\mathcal{E}_{n,\Lambda,2}(\Sigma)\geq 0$, completing the proof.
\end{proof}

\end{document}
